\section{Problem Formulation: Zero-Trust Memory Game}
\label{sec:problem}

This section formalizes memory-based multi-agent debate as a \textbf{zero-trust memory game}, where no agent's contribution to memory $\mathcal{M}$ is assumed honest.  Agents iterate between two roles: a \emph{contributor} who writes reasoning traces with factual claims into $\mathcal{M}$, and \emph{auditors} who inspect those contributions for hallucinations or adversarial content. Our goal is to design a set of game-theoretic rules, namely, the \textbf{utility functions} for contributor and auditors, respectively, to remove the need to trust any individual agent and achieve honest memory behavior on the whole-system.



\subsection{Setup and Notation for Shared-Memory Multi-Agent Debate}

We consider a shared-memory multi-agent debate system $\mathcal{G} = (\mathcal{N}, \mathcal{M})$ where $\mathcal{N} = \{a_1, \ldots, a_n\}$ is a set of $n$ agents operating over a
shared memory $\mathcal{M}$. We denote $\mathcal{M}_t$ as a memory state at each debate  round $t$. Besides agents' original roles, we assign two real-time roles at each round $t$: The agent that proposes new reasoning, facts, or beliefs is referred to as the \textbf{contributor}, denoted by $a_c \in \mathcal{N}$, who generates a proposal $\Delta_t$ to be written into $\mathcal{M}_t$. The remaining agents $\mathcal{A}_t = \mathcal{N} \setminus \{a_c\}$ form the \textbf{auditors}, tasked with evaluating whether $\Delta_t$ should be committed. A successful commit transitions the memory state as: \(\mathcal{M}_{t+1} = \texttt{Merge}(\mathcal{M}_t,\ \Delta_t)\) with an architecture-specific memory integration function $\texttt{Merge}(\cdot, \cdot)$ and the raw memory space $\mathbb{D}$ for all possible $\Delta_t$.
    
% \label{eq:merge}
% \end{equation}




\subsection{The Zero-Trust Memory Game}


\textbf{Rationale.} Recent work has shown that LLM-based agents are susceptible to
sycophantic agreement~\citep{fanous2025syceval,sharma2023towards}, misrepresentation~\citep{perez2022red}, and belief
perseveration~\citep{raj2022measuring}, all undermining the reliability of
agents acting as honest contributors or auditors.
Classical fault-tolerance models (e.g., Byzantine fault
tolerance~\citep{castro1999practical,lamport2019byzantine}) address a bounded
fraction of faulty agents but assume the remainder are honest, which is untenable in dynamic agent debating.
We therefore adopt a \emph{zero-trust} formulation: no agent is assumed to act
in alignment with collective memory integrity, and the system must enforce
correctness structurally rather than by trusting any participant.
See App.~\ref{app:zt_rationale} for a full discussion.

\begin{definition}[\textbf{Zero-Trust}]
\label{assm:zerotrust}
No agent in $\mathcal{N}$ is assumed to act in alignment with memory
integrity. Specifically: \textbf{(i) The contributor $a_c$} may craft $\Delta_t$ to advance a
    private objective $u_c$ misaligned with the true information state.\textbf{ (ii) Each auditor $a_i \in \mathcal{A}_t$} may act to advance its own private
    objective $u_i$, which may include collusion with $a_c$, or
    adversarial rejection of truthful updates.
\end{definition}

\noindent
This assumption \textbf{does not} require that all agents are adversarial, only
that none can be \emph{a priori} trusted.

\textbf{Game Formulation.} We formalize memory updating as a two-stage game: In the \textbf{contribution stage,} the contributor $a_c$ observes the current memory state $\mathcal{M}_t$ and a private
signal $\sigma_c \in \Sigma$ (representing its local reasoning context), and selects
a delta:
\begin{equation}
    \Delta_t \in \arg\max_{\Delta \in \mathbb{D}}\ u_c(\Delta,\ \mathcal{M}_t,\ \sigma_c)
\end{equation}
The private utility $u_c$ may encode objectives such as promoting a particular
conclusion, inflating confidence estimates, or injecting subtly biased representations
into $\mathcal{M}$. 

In the \textbf{auditing stage}, each auditor $a_i \in \mathcal{A}_t$ observes $\Delta_t$ and the current memory
$\mathcal{M}_t$, and independently emits a binary signal:
\begin{equation}
    v_i \in \{0, 1\}, \quad v_i = \pi_i(\Delta_t,\ \mathcal{M}_t,\ \sigma_i)
\end{equation}
where $\sigma_i \in \Sigma$ is $a_i$'s private signal and $\pi_i$ is its
(potentially mixed) strategy. The aggregate audit decision is formed by a
\textbf{commitment rule} $\mathcal{C}: \{0,1\}^{|\mathcal{A}_t|} \to \{0,1\}$,
such as a weighted majority vote. The update is committed if and only if
$\mathcal{C}(\{v_i\}_{i \in \mathcal{A}_t}) = 1$.

\subsection{Game-Theoretical Utility Functions and System-Level Equilibrium}
\label{ssec:equilibrium}

Each participant in the game defines a \textbf{utility function} to quantify
their payoffs in the memory game. At the system level, we leverage an
equilibrium of these payoffs to achieve optimal trust in memory. 

\textbf{Contributor utility.}
The contributor $a_c$ trades off the persuasiveness of $\Delta_t$
(the probability it is committed) against the risk of detection:
\begin{equation}
\label{eq:utility-contributor}
    u_c(\Delta_t,\ \mathcal{M}_t) =
        \underbrace{\Pr[\mathcal{C}(\{v_i\}) = 1 \mid \Delta_t]}_{\text{commitment gain}}
        \cdot \underbrace{\mathbb{I}(\Delta_t,\ \mathcal{M}_t)}_{\text{influence}}
        - \underbrace{D(\Delta_t,\ \mathcal{M}_t)}_{\text{detection penalty}}
\end{equation}
where $\Pr[\mathcal{C}(\{v_i\}) = 1 \mid \Delta_t]$ is the probability of
successfully committing $\Delta_t$, $\mathbb{I}(\Delta_t, \mathcal{M}_t)$
measures the downstream influence of $\Delta_t$ on future retrievals and
reasoning steps, and $D(\Delta_t, \mathcal{M}_t) \in [0,1]$ is the
detectability of $\Delta_t$ as inconsistent with $\mathcal{M}_t$.

\textbf{Auditor utility.}
Each auditor $a_i$
evaluates the \emph{alignment} of the proposed $\Delta_t$ with the
current memory state $\mathcal{M}_t$:
\begin{equation}
\label{eq:utility-auditor}
    u_i(v_i,\ \Delta_t,\ \mathcal{M}_t) =
        \underbrace{\texttt{Align}(\Delta_t,\ \mathcal{M}_t)}_{\text{consistency measure}}
        - \underbrace{c_i \cdot \mathbf{1}[v_i = \text{scrutinize}]}_{\text{effort cost}}
\end{equation}
where $\texttt{Align}(\Delta_t, \mathcal{M}_t) \in [-1, +1]$ measures the
semantic and logical consistency of $\Delta_t$ with existing memory, which is
positive when $\Delta_t$ corroborates or extends $\mathcal{M}_t$, negative
when it contradicts or destabilizes it. The $c_i > 0$ is the cost of
careful scrutiny. The alignment signal $\texttt{Align}(\cdot)$ is not
self-reported by $a_i$ but is computed by the calibration mechanism $\Phi$
from $a_i$'s submitted evidence (probe queries or attestation paths),
ensuring no auditor can manipulate their own utility measurement.
The effort cost $c_i$ is estimated online as the empirical computational
overhead of $a_i$'s scrutiny actions across rounds, requiring no
pre-specification (see App.~\ref{app:utility_rationale}).

\textbf{System-level equilibrium as the indicator of optimal trust in memory.}
With both utility functions defined above, the system reaches an equilibrium
when no participant can improve their payoff by unilaterally changing
their strategy: the contributor cannot craft a more influential $\Delta_t$
without increasing its detection risk, and no auditor can reduce their
scrutiny effort without degrading the alignment signal they produce.
Formally:
\begin{align}
\label{eq:equilibrium}
    \pi_i^* \in \arg\max_{\pi_i}\
        u_i(\pi_i,\ \Delta_t^*,\ \mathcal{M}_t)
        \quad \forall\, a_i \in \mathcal{A}_t \quad \texttt{and} \quad \Delta_t^* \in \arg\max_{\Delta \in \mathbb{D}}\
        u_c(\Delta,\ \mathcal{M}_t)
\end{align}
However, not all equilibria are desirable. Under multi-agent debating, three
integrity-degrading equilibria generically emerge: auditors
\emph{free-ride} by approving without scrutiny to avoid effort cost
$c_i$; \textit{colluding} auditors approve $\Delta_t$ unconditionally for
side benefits; or \textit{adversarial} auditors reject all updates regardless
of their quality (Propositions~\ref{prop:freeriding}--\ref{prop:adversarial_rejection},
App.~\ref{app:failure_modes}). These failures echo to incorrect memory commitment (Figure \ref{fig:intro}, detailed in App.~\ref{app:failure_modes}) and share a common cause:
the commitment outcome depends on auditor votes, giving agents wrong
incentives.

To avoid those failure modes, we develop calibration mechanism $\Phi$ to break such dependency, i.e.,
commitment is determined structurally by a calibration indicator $\rho$, not
by votes, which calibrate the memory system without requiring
any agent to be honest. Guided by a unified design rationale, we adapt $\Phi$ based on specific memory architecture (embedding-based or graph-based) in \S\ref{sec:calibration}. Due to space constraints, we defer the theoretical guarantees on equilibrium existence, as well as how calibration leads to optimal trust in memory, to App. \ref{app:equilibrium}.